\documentclass[10pt,a4paper,ragged2e,withhyper]{style} \geometry{left=1.25cm,right=1.25cm,top=1cm,bottom=1cm,columnsep=1.2cm} \usepackage{paracol} \usepackage{xcolor} \definecolor{MainGreen}{HTML}{1f2c67} \definecolor{SlateGrey}{HTML}{2E2E2E} \definecolor{LightGrey}{HTML}{666666}  \colorlet{heading}{MainGreen} \colorlet{headingrule}{MainGreen} \colorlet{accent}{MainGreen} \colorlet{emphasis}{SlateGrey} \colorlet{body}{LightGrey} \renewcommand{\cvItemMarker}{{\small\textbullet}} \renewcommand{\cvRatingMarker}{\faCircle} \renewcommand{\cvsectionfont}{\large\bfseries} \renewcommand{\cvsubsectionfont}{\normalsize\bfseries}  \renewcommand{\namefont}{\Huge\rmfamily\bfseries}
% --- Personnalisations de mise en forme propres à ce document (n'affectent
% pas style.cls ni les autres CV) : ligne décorative fine sous les titres de
% section, dates en italique, espacements resserrés pour un rendu plus
% compact et professionnel.
\renewcommand{\cvsection}[2][]{%
  \par\vspace{4pt}%
  \ifstrequal{#1}{}{}{\marginpar{\vspace*{\dimexpr1pt-\baselineskip}\raggedright\input{#1}}}%
  {\color{heading}\cvsectionfont \MakeUppercase{#2}}\par%
  \vspace{1pt}
  {\color{headingrule}\hdashrule{\linewidth}{0.5pt}{2pt}\par}\vspace{2pt}
}
\renewcommand{\cvevent}[4]{%
  \normalsize\par\noindent
  {\bfseries\color{emphasis}#1\ifstrequal{#2}{}{}{ -- #2}}%
  \ifstrequal{#3}{}{}{%
    \hfill{\color{accent}\bfseries\itshape\cvDateMarker~%
      \BeginAccSupp{method=pdfstringdef,ActualText={\datename: }}#3\EndAccSupp{}}%
  }\par
  \ifstrequal{#4}{}{}{%
    {\small\itshape\color{body}\cvLocationMarker~%
      \BeginAccSupp{method=pdfstringdef,ActualText={\locationname: }}#4\EndAccSupp{}\par}}%
  \vspace{2pt}%
  \footnotesize
}
\renewcommand{\divider}{\par\noindent\textcolor{body!30}{\hdashrule{\linewidth}{0.5pt}{0.4ex}}\vspace{2pt}}
\setlist{leftmargin=*,labelsep=0.5em,nosep,itemsep=0.08\baselineskip,after=\vspace{0.08\baselineskip}}
\begin{document} \name{El Moustpha Cheikh JIDDOU} \tagline{\textcolor{MainGreen}{Élève ingénieur à l'École Supérieure Polytechnique}} \photoR{2.5cm}{my photo}  \personalinfo{ \email{elmoustapha.cheikh.jiddou@gmail.com} \phone{+22248647376} \location{Nouakchott, Mauritanie} \printinfo{\faLinkedin}{El Moustpha Cheikh Jiddou}[https://www.linkedin.com/in/el-moustapha-cheikh-jidou/] \printinfo{\faGithub}{MoustaphaCheikhJidou}[https://github.com/MoustaphaCheikhJidou] \printinfo{\faKaggle}{jiddou26}[https://www.kaggle.com/jiddou26] \printinfo{\faGlobe}{Portfolio}[https://moustaphacheikhjidou.github.io/MyPortfolio/] } \makecvheader \columnratio{0.6} \begin{paracol}{2} \cvsection{Expériences Professionnelles}
\cvevent{Ingénieur Data \& IA}{Richat Partners}{Août 2026 -- Présent}{Nouakchott, Mauritanie}
\begin{itemize}
  \item Conception et implémentation de solutions IA et data pour les missions de conseil stratégique.
  \item Développement de pipelines de données et de modèles machine learning pour appuyer la prise de décision dans les projets de transformation digitale, d'investissement et de stratégie.
  \item Collaboration avec des équipes pluridisciplinaires pour traduire les besoins métiers en solutions techniques.
  \item Contribution aux offres IA du cabinet : intégration de LLM, systèmes RAG, analyse prédictive et conseil data-driven.
\end{itemize}
\divider
\cvevent{Stagiaire Architecte Systèmes de Données \& IT}{Banque Nationale de Mauritanie (BNM)}{Février 2025 -- Présent}{Nouakchott, Mauritanie}
\begin{itemize}
    \item Conception et implémentation d'un Data Warehouse centralisé d'entreprise s'appuyant sur des pipelines ELT robustes, l'orchestration automatisée des flux et des modèles prédictifs.
    
    \item Architecture et déploiement d'une infrastructure de microservices conteneurisée via des réseaux Docker isolés sur des serveurs applicatifs de production on-premise.
    
    \item Développement et intégration d'un agent conversationnel intelligent sécurisé (Click) connecté aux API REST de WhatsApp, incluant un module NLU personnalisé pour le routage dynamique des requêtes.
    
    \item Mise en œuvre d'une architecture RAG (Retrieval-Augmented Generation) avec base de données PostgreSQL/pgvector pour indexer et interroger la documentation bancaire officielle sans fuite de données.
    
    \item Évaluation technique et benchmarking de plusieurs grands modèles de langage (GPT-4o, Claude 3.5, Gemini 1.5, Llama local) afin d'optimiser le compromis performance, coûts d'infrastructure et stricte confidentialité.
    
    \item Développement de modèles de réseaux de neurones séquentiels (LSTM, GRU) infra-horaires pour anticiper les chutes rapides de puissance photovoltaïque et planifier la résilience des infrastructures critiques.
    
    \item \textbf{Environnement technique :} Python (FastAPI), Data Warehousing, pipelines ELT, Docker, N8N, PostgreSQL (pgvector), PyTorch, React, API REST, Linux/Bash.
\end{itemize}
\divider
\cvevent{Stage de fin d’études — Analyse macroéconomique et modélisation DSGE}{Banque Centrale de Mauritanie (BCM) — Direction Générale des Études et de la Stabilité Monétaire}{Février 2026 -- Mai 2026}{Nouakchott, Mauritanie}

\begin{itemize}
  \item Réalisation d’un mémoire de master en Économétrie et Statistique Appliquée portant sur l’analyse comparative des régimes monétaires en Mauritanie à l’aide d’un modèle DSGE de petite économie ouverte.
  
  \item Analyse des mécanismes de transmission de la politique monétaire, de l’inflation, du taux de change et des dynamiques macroéconomiques dans le contexte mauritanien.
  
  \item Exploitation et structuration de séries macroéconomiques et monétaires issues de la BCM, de l’ANSADE et de bases internationales pour la calibration et les simulations économiques.
  
  \item Participation aux activités de la Direction Générale des Études et de la Stabilité Monétaire : analyse conjoncturelle, prévisions macroéconomiques, statistiques monétaires et études économiques.
  
  \item Développement d’outils de modélisation et de simulation économique pour comparer plusieurs régimes de politique monétaire et évaluer leur impact sur la stabilité macroéconomique.
  
  \item Production d’analyses quantitatives, graphiques économiques et indicateurs statistiques pour l’interprétation des résultats et l’aide à la décision.
  
  \item \textbf{Environnement :} Econométrie, DSGE, analyse macroéconomique, statistiques monétaires, Python, R, Excel, Julia, GAMS, modélisation économique, analyse de données.
\end{itemize}

\divider
\cvevent{Enseignant vacataire — Pauvreté \& Conditions de Vie des Ménages + Enquête Sociologique}{Institut Supérieur des Métiers de la Statistique (ISS)}{Février 2026 -- Juin 2026}{Nouakchott, Mauritanie}
\begin{itemize}
  \item Reconduction pour une deuxième année avec un module supplémentaire.
  \item Enseignement du module Pauvreté et Conditions de Vie des Ménages (édition améliorée).
  \item Conception et enseignement d'un nouveau module : Enquête Sociologique (méthodologie d'enquête, échantillonnage, conception de questionnaires, analyse des données).
  \item Encadrement des TD et TP avec simulation d'enquêtes terrain et utilisation de KoboToolbox, CSPro et Stata.
  \item Direction et évaluation des PFE et encadrement des stages ouvriers.
  \item Développement des supports de cours, études de cas et grilles d'évaluation.
\end{itemize}
\divider
\cvevent{Enseignant vacataire — Pauvreté et Conditions de Vie des Ménages}{Institut Supérieur des Métiers de la Statistique (ISS)}{Février 2025 -- Juin 2025}{Nouakchott, Mauritanie}
\begin{itemize}
  \item Première expérience d'enseignement supérieur au sein de l'établissement dont je suis diplômé.
  \item Préparation et animation des cours magistraux sur les méthodologies de mesure de la pauvreté et l'analyse des conditions de vie des ménages.
  \item Conception et encadrement des travaux dirigés (TD) et travaux pratiques (TP), avec études de cas réels issus d'enquêtes ménages nationales.
  \item Accompagnement pédagogique personnalisé et suivi continu des étudiants tout au long du semestre.
  \item Encadrement des projets de fin d'études (PFE) et des stages ouvriers, avec mentorat sur les aspects méthodologiques et techniques.
\end{itemize}
\divider
\cvevent{Stagiaire assistant-ingénieur — Département Mine, Service Géotechnique}{Kinross Gold Corporation (Tasiast Mauritanie)}{Juillet 2025 -- Septembre 2025}{Tasiast, Mauritanie}
\begin{itemize}
  \item \textbf{Wall Compliance \& 80/80} : conception d’un mini-dashboard Python/Gradio pour le suivi \textit{Toe/Crest/Berm Width/BFA} à partir des fichiers \texttt{CALC} et \texttt{COMPLIANCE} (auto-détection des entrées, KPI, barres 80/80, tendances \& deltas).
  \item \textbf{Digitalisation Daily Survey} : chaîne \textbf{Microsoft} \textit{Forms} $\rightarrow$ \textit{SharePoint/Excel Web} $\rightarrow$ \textit{Power BI} avec typages/contrôles, calcul de durées, KPI (\% Done, Backlog, SLA Breach) et rafraîchissement quasi temps réel.
  \item \textbf{Prévision PV infra-horaire} : préparation de données minute (\textit{lags}, rampes \(\Delta P\), encodages temporels) et prototypes séquentiels (LSTM/GRU) pour anticiper les chutes rapides de puissance et améliorer le \textit{lead time} d’alerte.
  \item \textbf{Pratiques} : versionnement \texttt{git}, documentation technique, transfert de compétence aux équipes terrain/process.
  \item \textbf{Environnement} : Python (NumPy, pandas, Plotly, Gradio), Power BI, SharePoint/Excel Web.
\end{itemize}
\divider

\cvevent{Président \& Coordinateur d'événements}{ESP Data Club}{Octobre 2022 -- Avril 2025}{Nouakchott, Mauritanie} \begin{itemize} \item \textbf{Coordinateur d'événements} (Octobre 2022 -- Janvier 2023) : Organisation d'événements pour promouvoir la statistique et l’analyse de données. Coordination avec les membres pour assurer le succès des activités. \item \textbf{Président} (Janvier 2023 -- Avril 2025) : Direction du club, établissement des objectifs stratégiques, gestion de projets de collecte et d'analyse de données, développement de partenariats académiques. \item Compétences : Leadership, gestion de projet, planification \\ d'événements, communication. \item\href{https://esp-data-club.com/}{Voir le Club}. \end{itemize} \divider \cvevent{Stage De Fin D’étude}{Caisse des Dépôts et de Développement (CDD)}{Février 2023 -- Juin 2023}{Nouakchott, Mauritanie} \begin{itemize} \item Analyse des données bancaires issues de la CDD pour évaluer l'impact du système de garantie de crédit sur les défauts de paiement en Mauritanie. \item Utilisation de modèles économétriques et d'approches modernes de machine learning pour prédire les risques de défaut. \item Élaboration et interprétation de rapports statistiques concernant les caractéristiques sociodémographiques et financières des emprunteurs. \item Développement de solutions computationnelles pour l'analyse des risques avec des algorithmes tels que la régression logistique, k-NN et Random Forest. \item\href{https://github.com/MoustaphaCheikhJidou/credit-default-analysis.git}{Voir les analyses}. \end{itemize} \divider \cvevent{Agent Recenseur}{Agence Nationale de la Statistique et de l'Analyse Démographique et Économique (ANSADE)}{Juillet 2022 -- Août 2022}{Nouakchott, Mauritanie} \begin{itemize} \item Collecte de données sur le terrain pour le RGPH pilote. \item Compétences acquises : collecte de données, communication, travail sur le terrain. \end{itemize} \divider \cvevent{Stage Ouvrier}{Agence Nationale de la Statistique et de l'Analyse Démographique et Économique (ANSADE)}{Avril 2022 -- Juin 2022}{Nouakchott, Mauritanie} \begin{itemize} \item Élaboration d'une méthodologie pour la cartographie de données en vue de la conception des comptes sectoriels. \item Participation à la collecte et au traitement des données des branches économiques pour les comptes nationaux. \item Analyse des données sectorielles en utilisant le logiciel ERETES pour compiler les comptes économiques. \item Formation sur les techniques de décomposition du PIB dans le cadre du Programme de Comparaison Internationale. \item Contribution au processus d'enquête nationale pour la mobilisation et la quantification de la ZAKAT, incluant la formation des enquêteurs. \end{itemize} 

\cvsection{Projets}
\cvevent{Grain Size Analysis of 316L Microstructure}{Projet académique : Analyse granulométrique par vision par ordinateur}{Juin 2025 -- Août 2025}{Nouakchott, Mauritanie}
\begin{itemize}
  \item Développement d’un pipeline complet de traitement d’images pour segmenter les joints de grains d’un acier inox 316L à partir de micrographies à 500×.
  \item Mise en œuvre et comparaison de méthodes de segmentation : seuillage classique, filtres de gradient, modèle HED (Caffe/OpenCV) et U-Net convolutif entraîné sur données réelles et artificielles.
  \item Extraction automatique de la granulométrie en unités physiques (µm, µm²) via calibration pixels/micron et calcul de statistiques (distribution des tailles, axes, circularité).
  \item Création d’une interface interactive (Gradio) permettant de charger une micrographie, segmenter les grains et exporter les caractéristiques géométriques en CSV.
  \item \href{https://github.com/MoustaphaCheikhJidou/grain-size-analysis-316L-microstructure}{Voir le dépôt GitHub}.
\end{itemize}
\divider


\cvevent{RAG‑WebSystem}{Projet académique : Système d’IA générative à base de Retrieval‑Augmented Generation}{Mars 2025 -- Avril 2025}{Nouakchott, Mauritanie} \begin{itemize} \item Conception d’un moteur RAG (Retrieval + LLM) capable de répondre à des requêtes en s’appuyant sur des documents externes indexés par vecteurs. \item Implémentation d’une chaîne complète : ingestion, vectorisation (FAISS), recherche sémantique, puis génération de réponse contextuelle avec un modèle open‑source. \item Construction d’une API REST (FastAPI) et d’une interface web réactive (React + Tailwind) pour tester, tracer et évaluer les réponses. \item\href{https://github.com/MoustaphaCheikhJidou/RAG-WebSystem}{Voir le dépôt GitHub}. \end{itemize} \divider \cvevent{DataMine App}{Projet académique : Projet industriel d'entreprise}{Novembre 2023 -- Aujourd’hui}{Nouakchott, Mauritanie} \begin{itemize} \item Conception et développement d'une plateforme collaborative pour la gestion et le suivi des tâches d'un projet de données géostatistiques. \item Analyse des données géostatistiques pour optimiser la gestion des ressources et améliorer la prise de décision. \item Création d'une interface utilisateur intuitive pour la gestion et visualisation des tâches sur Github \item\href{https://github.com/MoustaphaCheikhJidou/DataMineSynergy-App.git}{Voir le projet ici}. \end{itemize} \divider \cvevent{Système d’Analyse de Données Agricoles (Serre Agricole)}{Projet académique : Projet entrepreneuriat et innovation}{Novembre 2023 -- Juin 2024}{Nouakchott, Mauritanie} \begin{itemize} \item Conception et développement d’un système intelligent pour l’analyse de données agricoles, spécifiquement pour les serres agricoles. \item Implémentation de capteurs IoT pour la collecte des données en temps réel et création de tableaux de bord interactifs pour le suivi. \item Application de techniques de machine learning pour optimiser la gestion des ressources et améliorer les rendements agricoles. \item\href{https://github.com/MoustaphaCheikhJidou/smart-agriculture-system.git}{Voir le projet ici}. \end{itemize} \divider \cvevent{Classification NLP des Tweets en Cas de Désastre}{Projet académique : Introduction aux big data}{Janvier 2022 -- Février 2022}{Nouakchott, Mauritanie} \begin{itemize} \item Développement d’un modèle d’apprentissage automatique pour classifier les tweets relatifs aux désastres naturels. \item Traitement de données massives et application de techniques avancées de NLP (Natural Language Processing). \item Utilisation de bibliothèques Python pour la construction et l'évaluation du modèle. \item\href{https://github.com/MoustaphaCheikhJidou/disaster-tweets-nlp-classification.git}{Voir le projet ici}. \end{itemize} \divider \cvevent{Application de Gestion d’Enquête}{Projet académique : Base de données avancées}{Décembre 2022 -- Janvier 2022}{Nouakchott, Mauritanie} \begin{itemize} \item Conception et réalisation d’une application de gestion d'enquête pour faciliter la collecte et le suivi des données. \item Développement backend avec le framework Django et création du front-end en HTML, CSS et JavaScript. \item\href{https://github.com/MoustaphaCheikhJidou/survey-data-collection.git}{Voir le projet ici}. \end{itemize} \divider \cvevent{Enquête sur la perception des Influenceurs}{Projet académique : Enquêtes sociologiques}{Janvier 2022 -- Mars 2022}{Nouakchott, Mauritanie} \begin{itemize} \item Conception et réalisation d’une enquête sur la perception des influenceurs dans la société mauritanienne. \item Élaboration du questionnaire sur KoboCollect pour assurer une collecte de données structurée. \item Collecte et analyse autonome des données, avec interprétation des résultats pour dégager des tendances sociales. \end{itemize} \divider \cvevent{Application de Gestion de Notes}{Projet académique : Base de données et algorithmes avancés}{Décembre 2022 - Janvier 2022}{Nouakchott, Mauritanie} \begin{itemize} \item Conception et réalisation d’une application de gestion de notes pour les étudiants. \item Développement backend avec le framework Flask et utilisation des langages HTML, CSS et JavaScript pour le front-end. \item\href{https://github.com/MoustaphaCheikhJidou/ISMSNotesApp.git}{Voir le projet ici}. \end{itemize} \divider \cvevent{Les Déterminants des Performances Scolaires en Mauritanie}{Projet académique : Logiciel Statistique}{Mai 2021 -- Juin 2021}{Nouakchott, Mauritanie} \begin{itemize} \item Analyser les facteurs influençant les performances scolaires en Mauritanie sous R. \item Étudié les données issues des systèmes éducatifs de la CONFEMEN (PASEC1), comprenant des variables socio-démographiques, géographiques, et pédagogiques. \end{itemize} \cvsection{Savoir-Être}
{\small Force de proposition \textbullet{} Curiosité \textbullet{} Aisance relationnelle \textbullet{} Créatif \textbullet{} Capacité d'écoute}

\switchcolumn \cvsection{Formation}
\cvevent{École Supérieure Polytechnique de Nouakchott}{ESP - Mauritanie}{Octobre 2023 -- Juin 2026}{} Diplôme : Ingénieur d'État en Statistique et en Ingénierie des données.

\divider

\cvevent{École Nationale Supérieure d'Informatique et d'Analyse des Systèmes}{ENSIAS - Rabat}{Octobre 2025 -- Février 2026}{} Semestre d’échange : Ingénierie Informatique et Réseaux.

\divider

\cvevent{Faculté des Sciences Économiques et de Gestion}{FSEG - Mauritanie}{Octobre 2023 -- Juin 2025}{} Diplôme : Master en Économétrie et Statistiques Appliquées -- Parcours Ingénierie Statistique Économique.
\divider

\cvevent{Groupe Polytechnique-Institut Supérieur des Métiers de la Statistique}{GP - Mauritanie}{Octobre 2020 -- Juin 2023}{} Diplôme : Licence Professionnel en Statistique Appliquée.\\ \divider \cvevent{Lycée d'Atar}{Atar – Mauritanie}{2019 -- 2020}{} Diplôme : Baccalauréat option Mathématiques (Bac C).

\cvsection{Compétences}
{\small
\begin{itemize}
\item \textbf{Analyse \& modélisation} : statistique, économétrie, sondages \& enquêtes; data mining (classification, régression, clustering, règles d’association); séries temporelles (prévision courte portée).
\item \textbf{Big Data, Data Engineering, DB \& SIG} : Architecture \& modélisation de Data Warehouse centralisé, pipelines ELT/ETL (Airbyte, dbt), orchestration de flux de données (N8N), Hadoop, Spark, Kafka, MapReduce, Hive SQL; gestion \& optimisation SQL (plans d’exécution, indexation, partitionnement, Oracle SQL Tuning); Docker \& Dev Containers; Stata, SPSS, Excel, GAMS, Eviews, Sphinx, KoboToolbox, ODKCollect, CSPro; Oracle SQL Developer; PostgreSQL/PostGIS (incl. pgvector); QGIS/ArcGIS; intégration \& analyse SIG (jointures spatiales, géotraitements de base).
\item \textbf{IA, MLOps \& programmation} : ML supervisé (régression, SVM, forêts aléatoires, arbres de décision); deep learning (LSTM/GRU pour séries temporelles, CNN, Transformers) avec TensorFlow, PyTorch, Keras; NLP \& RAG (embeddings, modèles de langage, recherche augmentée, RAGAI API); MLOps (MLflow, Docker/Compose); Python/PySpark, Julia, R, Scala, SQL (incl. HiveQL), PL/SQL, Java (J2EE, \textbf{Spring Boot, Spring Cloud, Spring Data, Spring Security, microservices Spring}), VBA, \LaTeX; Bash \& Linux CLI; Git/GitHub; Jira, Trello, VS Code, Cursor; développement web (HTML, CSS, JavaScript, Django, Flask, APIs REST, intégration middleware J2EE), React, Flutter.
\end{itemize}
}

\cvsection{Langues} \cvskill{Arabe}{4}{Langue natale}\\ \divider \cvskill{Français}{3}{DALF C1}\\ \divider \cvskill{Anglais}{2} {CEFR B1 - Tracktest English Assessment}

\cvsection{Vie Associative} \begin{itemize} \item Président de l'Alumni de l'ISS (poste actuel) \item Responsable des événements et programmes, Association des Statisticiens Mauritaniens \item Président de l'ESP Data Club (Janvier 2023 -- Avril 2025) \item Coordinateur d'événements à l'ESP Data Club (Octobre 2022 - Janvier 2023) \item Membre du Clubs Leaders Étudiants Francophones (CLEF) \item Membre Mauritanian Engineering Community (MEC) \item Formation Militaire \end{itemize} \cvsection{Fiertés} \cvachievement{\faTrophy}{1ère place au RIMAI Data Science Challenge 2024}{Compétition axée sur la reconnaissance de plaques d'immatriculation en utilisant des techniques d'IA. En savoir plus sur le défi à \href{https://www.rim-ai.com/challenge}{RIMAI Challenge Website}.} 
 
 \divider 
 
 \cvachievement{\faUsers}{Organisation du premier Datathon en Mauritanie}{Compétition de trois jours axée sur l'analyse des données mauritaniennes. En tant que président du ESP Data Club, j'ai supervisé l'organisation et évalué les challengeurs à travers des défis compétitifs en data science.}
 
 \cvsection{Certifications}
{\small
\textbf{freeCodeCamp (6) :} \href{https://freecodecamp.org/certification/el_moustapha_cheikh_jidou/relational-database-v8}{Relational Database (v8)} (2025) \textbullet{} \href{https://freecodecamp.org/certification/EL_Moustapha_Cheikh_Jidou/data-visualization}{Data Visualization} (2023) \textbullet{} \href{https://freecodecamp.org/certification/EL_Moustapha_Cheikh_Jidou/data-analysis-with-python-v7}{Data Analysis with Python} (2022) \textbullet{} \href{https://freecodecamp.org/certification/EL_Moustapha_Cheikh_Jidou/scientific-computing-with-python-v7}{Scientific Computing with Python} (2022) \textbullet{} \href{https://freecodecamp.org/certification/EL_Moustapha_Cheikh_Jidou/front-end-development-libraries}{Front End Development Libraries} (2022) \textbullet{} \href{https://freecodecamp.org/certification/ELMoustapha_Cheikh_Jidou/responsive-web-design}{Responsive Web Design} (2022)\\[3pt]
\textbf{DataCamp (9) :} \href{https://www.datacamp.com/statement-of-accomplishment/course/0b455b5c403a108e847df406cae8ae46aa9623f6?raw=1}{Data Manipulation in SQL} (2023) \textbullet{} \href{https://www.datacamp.com/statement-of-accomplishment/course/dc4900c31a4d547e6e4f317fe9f64588c49fb71e?raw=1}{Intermediate SQL} (2023) \textbullet{} \href{https://www.datacamp.com/statement-of-accomplishment/course/b1207834324f0c856760624c2bb470b3afd6d642?raw=1}{Introduction to Data Science in Python} (2023) \textbullet{} \href{https://www.datacamp.com/statement-of-accomplishment/course/c6351222bc9ed2c5771273e520045c09e60d5b48?raw=1}{Python Data Science Toolbox (Part 2)} (2023) \textbullet{} \href{https://www.datacamp.com/statement-of-accomplishment/course/c129bc893ba40d2669b9682325ad8a4c3b6d8ecc?raw=1}{Understanding Data Engineering} (2023) \textbullet{} \href{https://www.datacamp.com/statement-of-accomplishment/course/6995ad305b5768d3d465ca0c38f3c26b8b25e3a4?raw=1}{Writing Efficient R Code} (2023) \textbullet{} \href{https://www.datacamp.com/statement-of-accomplishment/course/43feb083d3979dbbec34f54a6c8cf6458e15cd12}{Intermediate Python} (2022) \textbullet{} \href{https://www.datacamp.com/statement-of-accomplishment/course/929a5e8c2a4e42f4880ba65107b27e8f4b411fa3}{Intermediate R} (2022) \textbullet{} \href{https://www.datacamp.com/statement-of-accomplishment/course/a0d86c15550db6e4e8955c19e2c1f8037bda5497}{Introduction to Version Control with Git} (2022)\\[3pt]
\textbf{Google (3) :} \href{https://www.coursera.org/account/accomplishments/records/UCU4W3HV5KAR}{Google Advanced Data Analytics Capstone} (2023) \textbullet{} \href{https://www.coursera.org/account/accomplishments/specialization/F6NY5HAVFY2S}{Google Advanced Data Analytics Certificat Professionnel} (2023) \textbullet{} \href{https://www.coursera.org/account/accomplishments/specialization/certificate/JG3NPUJ7LZGL}{Google Data Analytics} (2023)
}

 \cvsection{Centres d'Intérêts} \cvtag{Coding}\cvtag{Voyage}\cvtag{Football}\cvtag{Tourisme} \end{paracol} 
\end{document}